\section{INTRODUCTION}
\label{sec:introduction}
% =============================================================================

Ngành thép chiếm khoảng 7--9\% lượng phát thải CO$_2$ nhân tạo toàn cầu, với tiêu thụ điện năng đại diện cho một chi phí vận hành chủ đạo. Do đó, dự báo tải chính xác và phát hiện sớm các bất thường điện là rất quan trọng cho tối ưu hóa năng lượng và bảo trì dự đoán trong các nhà máy sản xuất thép.

Mặc dù có sẵn các tập dữ liệu năng lượng công nghiệp công khai, các kho ngữ liệu hiện có thường chỉ cung cấp các bản ghi đồng hồ thô mà không có các chỉ báo vật lý phái sinh hoặc chú thích bất thường. Hạn chế này làm giảm khả năng áp dụng của các mô hình học có giám sát cho chẩn đoán lỗi. Hơn nữa, các sự kiện bất thường---chẳng hạn như động cơ chạy không tải, suy giảm cách điện, và quá tải cục bộ---là hiếm, dẫn đến mất cân bằng lớp nghiêm trọng làm giảm hiệu suất bộ phân loại.

Trong thực tế vận hành, các nhà máy có thể đã có relay bảo vệ, cảnh báo SCADA, hoặc hệ thống condition monitoring chuyên dụng để xử lý sự cố điện cấp an toàn. Tuy nhiên, các hệ thống này thường phụ thuộc vào cảm biến tài sản, log bảo trì, hoặc dữ liệu nội bộ khó công khai. Khoảng trống mà nghiên cứu này nhắm tới vì vậy không phải là thay thế hệ thống bảo vệ thời gian thực, mà là xây dựng một pipeline dữ liệu ít phụ thuộc hạ tầng cảm biến chuyên dụng: từ dữ liệu công tơ 15 phút có thể tái lập, bổ sung ngữ cảnh thời tiết, đặc trưng thời gian, wavelet, vật lý, và nhãn weak/proxy có thể kiểm toán cho phân tích offline.

Để giải quyết các hạn chế này, bài báo đề xuất một pipeline tiền xử lý và xây dựng tập dữ liệu có hệ thống cho dữ liệu điện ngành thép. Các đóng góp bao gồm ba phần: (i) một framework kỹ thuật đặc trưng đa miền bổ sung các tín hiệu thô với thống kê trễ thời gian, hệ số Biến đổi Wavelet rời rạc (DWT), và các đạo hàm mang ý nghĩa vật lý bao gồm Công suất biểu kiến và Góc lệch pha; (ii) một lược đồ gán nhãn bất thường dựa trên vật lý nhận diện các sự kiện chạy không tải, rò rỉ năng lượng, và quá tải cục bộ; và (iii) một Mạng đối kháng Sinh (GAN) để tăng cường tổng hợp lớp thiểu số nhằm giảm thiểu mất cân bằng dữ liệu.

Phần còn lại của bài báo được tổ chức như sau. Phần~\ref{sec:related} tổng quan các nghiên cứu liên quan. Phần~\ref{sec:methodology} trình bày chi tiết phương pháp đề xuất. Phần~\ref{sec:results} trình bày kết quả và thảo luận. Cuối cùng, Phần~\ref{sec:conclusion} kết luận bài báo.

% =============================================================================
